\begin{table}[H]
\centering
\caption{Summary Statistics by Sex and Age Group}
\begin{threeparttable}
\begin{tabular}{lcccccccc}
\toprule
 & \multicolumn{2}{c}{Employment} & \multicolumn{2}{c}{Sep.\ Rate} & \multicolumn{2}{c}{Hire Rate} & \multicolumn{2}{c}{Earnings (\$)} \\
\cmidrule(lr){2-3} \cmidrule(lr){4-5} \cmidrule(lr){6-7} \cmidrule(lr){8-9}
Group & Mean & SD & Mean & SD & Mean & SD & Mean & SD \\
\midrule
Female 25-34 & 244,887 & 276,400 & 0.2150 & 0.0426 & 0.2145 & 0.0457 & 2,802 & 726 \\
Female 45-54 & 231,391 & 250,442 & 0.1239 & 0.0261 & 0.1218 & 0.0269 & 3,624 & 1,076 \\
Male 25-34 & 269,122 & 306,403 & 0.2289 & 0.0428 & 0.2304 & 0.0461 & 3,806 & 875 \\
Male 45-54 & 254,605 & 281,794 & 0.1307 & 0.0284 & 0.1279 & 0.0286 & 6,149 & 1,655 \\
\bottomrule
\end{tabular}
\begin{tablenotes}[flushleft]
\small
\item Notes: N = 18,072 state-quarter-sex-age observations across 51 states, 2000--2022. Employment is beginning-of-quarter count. Separation rate = quarterly separations / employment. Hire rate = quarterly hires / employment. Earnings are average monthly earnings for stable workers. Source: Census Quarterly Workforce Indicators (QWI).
\end{tablenotes}
\end{threeparttable}
\label{tab:summary}
\end{table}
